Notation as in the statement; $k$ denotes the pair number $N_{P}$ of a basis state, so that $\sum_{i \le 3}\beta_{i} = k - \tfrac32$, $\mu_{P} = \tfrac18(k - \tfrac32)$ and $B - L = 1 - \tfrac{2k}{3}$ on that state (Proposition \claimref{synthesis.georgi-jarlskog-majorana-moment-map}(i)).

(i) For the Majorana element $T$ of Proposition \claimref{standard-model.majorana-element}, by Proposition \claimref{synthesis.georgi-jarlskog-square-of-majorana-element}(i), $TT^{\dagger} = \prod_{i}(\tfrac12 + \beta_{i})$ and its degree-four part is $G_{4} = \tfrac18\sum_{i < l}\beta_{i}\beta_{l}$, each $\beta_{i}\beta_{l}$ ($i \ne l$) a single monomial of degree four supported on the modes $i$ and $l$. The pair-only part is $P_{4} = \tfrac18(\beta_{1}\beta_{2} + \beta_{1}\beta_{3} + \beta_{2}\beta_{3})$. Since $\beta_{i}^{2} = \tfrac14$, $P_{4} = \tfrac{1}{16}\bigl((\beta_{1} + \beta_{2} + \beta_{3})^{2} - \tfrac34\bigr) = \tfrac{1}{16}(64\mu_{P}^{2} - \tfrac34) = 4\mu_{P}^{2} - \tfrac{3}{64}$. With $\mu_{P} = -\tfrac{3}{16}(B - L)$, $4\mu_{P}^{2} = \tfrac{9}{64}(B - L)^{2}$, which gives $P_{4} = \tfrac{3}{64}(3(B - L)^{2} - 1)$. On the four values of $k$ this is $\tfrac{3}{32}, -\tfrac{1}{32}, -\tfrac{1}{32}, \tfrac{3}{32}$, and $(B - L)(-1)^{k}$ takes the values $1, -\tfrac13, -\tfrac13, 1$, so $P_{4} = \tfrac{3}{32}(B - L)(-1)^{N_{P}}$. The pair-rung part is $X_{4} = 8\mu_{R}\mu_{P}$ (loc. cit. (ii)). The rung-only part is $R_{4} = \tfrac18\beta_{4}\beta_{5}$; since $(-1)^{n_{4} + n_{5}} = (-2\beta_{4})(-2\beta_{5}) = 4\beta_{4}\beta_{5}$, $R_{4} = \tfrac{1}{32}(-1)^{N_{R}} = -\tfrac{1}{32}\chi$, and $4\mu_{R}^{2} - \tfrac{1}{32} = \tfrac{1}{16}(\tfrac12 + 2\beta_{4}\beta_{5}) - \tfrac{1}{32} = \tfrac18\beta_{4}\beta_{5}$. $P_{4}$ is a function of the pair occupations symmetric in the three pairs, so it commutes with $\su(3)$ (which preserves $k$ and permutes the pair occupations within each $k$), with $B - L$, and with every rung operator, hence with $\su(2)_{L}$ and $\su(2)_{R}$; the bracket $v_{1}v_{2}$ raises $k$ by two, between values $\tfrac{3}{32}$ and $-\tfrac{1}{32}$ ($k = 0 \to 2$), so $P_{4}$ does not commute with it. Degree-four monomials in the six pair generators span $\Lambda^{4}$ of the pair space, which under $\su(4)$ is the adjoint $15$ (it is $\Lambda^{2}$ times the pair volume element); with the commutation just shown, $P_{4}$ lies in the $(15,1,1)$. $R_{4}$ is a multiple of $\chi$, the Pati-Salam singlet of Proposition \claimref{standard-model.automorphisms-preserve-slot-pair-split-are}.

(ii) The sixteen are the states of even total number (the particle half of Proposition \claimref{standard-model.hyp-s-slot-statistics}, with $\nu^{c}$ the empty state, the halving forced by the Pati-Salam stage, Theorem \claimref{synthesis.neutral-singlet-bit-decided-by-gauge-algebra-leaving-u5}(iii)), so $(-1)^{N_{P}} = (-1)^{N_{R}} = -\chi$ on it. For $\varphi$ in the sixteen, $w_{4}P_{4}\varphi = \tfrac{3}{32}w_{4}(B - L)(-1)^{N_{P}}\varphi = -\tfrac{3}{32}w_{4}(B - L)\chi\varphi = -\tfrac{3}{32}T_{\mathrm{GJ}}\varphi$. The mass forms on the sixteen involve $X\varphi$ for $\varphi$ in the sixteen only, so the form of $w_{4}P_{4}$ is $-\tfrac{3}{32}$ times that of $T_{\mathrm{GJ}}$, which by Proposition \claimref{families.doublet-direction-module-georgi-jarlskog-clebsch}(i), (ii) is symmetric, is the $126$ component of $16 \otimes 16$ (no $10$, no $120$), and has $|y_{e}/y_{d}| = 3$. The same argument with $v_{4}$ and part (iv) of that proposition gives the up-type statement.

(iii) By Theorem \claimref{standard-model.yukawa-pattern-two-rung-module}, with the pairing of Lemma \claimref{standard-model.invariant-pairing-five-modes}, $w_{4}$ couples exactly $(L, e^{c})$ and $(Q, d^{c})$ and its form is symmetric; every coupled pair is $B - L$ neutral (Proposition \claimref{families.doublet-direction-module-georgi-jarlskog-clebsch}(iii)), with $|B - L| = 1$ on the lepton pair and $\tfrac13$ on the quark pair (Proposition \claimref{standard-model.generation-ten-under-pati-salam}). Since $f(B - L)$ is diagonal, $B(\psi, w_{4}f(B - L)\varphi) = f(q_{\varphi})B(\psi, w_{4}\varphi)$, and exchanging $\psi$ and $\varphi$ replaces $q_{\varphi}$ by $-q_{\varphi}$. So the symmetric part is the $10$ form weighted by the even part $f_{e}$ of $f$ (by $f_{e}(1)$ on leptons, $f_{e}(\tfrac13)$ on quarks) and the antisymmetric part the same with the odd part $f_{o}$. By (ii) the Georgi-Jarlskog form is the $10$ form weighted by $-1$ on leptons and $\tfrac13$ on quarks. Solving $f_{e}(1) = x - y$, $f_{e}(\tfrac13) = x + \tfrac{y}{3}$ gives $x = \tfrac14(3f_{e}(\tfrac13) + f_{e}(1))$, a quarter of the trace of $f_{e}$ over a quartet, and $y = \tfrac34(f_{e}(\tfrac13) - f_{e}(1))$. For $f = B - L$ the antisymmetric part is by definition the grade-three form of Proposition \claimref{standard-model.grade-three-georgi-jarlskog-direction}. For the traces: a pair monomial that is not a product of whole mode bivectors $a_{i}b_{i} = -2\beta_{i}$ changes some occupation and has zero diagonal. On the even half $\{\emptyset, \{1,2\}, \{1,3\}, \{2,3\}\}$, $\beta_{i}$ is $+\tfrac12$ on the two states containing $i$ and $-\tfrac12$ on the other two, so its trace is $0$; $\beta_{i}\beta_{l}$ ($i \ne l$) is $+\tfrac14$ on $\emptyset$ and on $\{i, l\}$ and $-\tfrac14$ on the other two, trace $0$; and similarly on the odd half by the particle-hole map $k \mapsto 3 - k$, which sends $\beta_{i}$ to $-\beta_{i}$. The degree-six product $\beta_{1}\beta_{2}\beta_{3}$ is $-\tfrac18$ on every state of the even half, trace $-\tfrac12 \ne 0$. $P_{4}$ is even in $B - L$ with $3P_{4}(\tfrac13) + P_{4}(1) = -\tfrac{3}{32} + \tfrac{3}{32} = 0$, so $x = 0$, as (ii) states.

(iv) $X_{4}$: Proposition \claimref{synthesis.georgi-jarlskog-square-of-majorana-element}(iii). $R_{4} = -\tfrac{1}{32}\chi$, and on each coupled pair one state is a doublet ($\chi = 1$) and the other a singlet ($\chi = -1$), so $w_{4}\chi$ has antisymmetric form only. Adding, the symmetric parts are $-\tfrac{3}{32} + \tfrac{3}{32} = 0$ times the Georgi-Jarlskog form, plus nothing, so $w_{4}G_{4}$ has no symmetric form, as loc. cit. (iv) states: the $10$ relation with multiple zero. Finally $\tfrac{1}{32} + P_{4} = \tfrac{1}{64}(9(B - L)^{2} - 1)$ is $\tfrac18$ on the lepton values and $0$ on the quark values, which by (iii) weights the $10$ form by $\tfrac18$ on $(L, e^{c})$ and $0$ on $(Q, d^{c})$.

The check builds the module as $32 \times 32$ rational matrices, expands $TT^{\dagger}$ in all $1024$ Majorana monomials by the trace form, reassembles it, splits the degree-four part by support, and confirms (i) to (iv) as operator identities and as mass forms against the stored pairing, including the coefficients $(x, y)$ of (iii) for four even weightings and the quartet traces of every pair monomial of degrees two, four and six.
